\section{The remaining irregular-prime conditions}
\label{sec:irreg-gates}

The finite sets above do not by themselves identify the exact
centres in the cusp residue discs. We describe a sufficient
additional input for that local step, and then the distinct
restriction-map issue.

\subsection{A prescribed supported pairing}

In this subsection assume \(p>5\), fix a wild quotient \(C=C_{a,b,c}\),
and put \(A=\Jac(C)\), \(K=\Q(\zeta_p)\), and
\(\pi=1-\zeta_p\), acting on \(A/K\) by complex multiplication.
Let \(\mathfrak p\) be the unique prime of \(K\) above \(p\), and set
\begin{equation}
 \Lambda=
 \bigl(A(K_{\mathfrak p})/\pi A(K_{\mathfrak p})\bigr)^\Delta.
 \label{eq:irreg-Lambda}
\end{equation}
The local calculation in
\cite[Corollary 2 and Proposition 5, pp.~574--575]{McCallum1994}
shows that \(\Lambda\) is one-dimensional over \(\F_p\), and that
the natural map
\begin{equation}
 \chi:A(\Q_p)/pA(\Q_p)\longrightarrow\Lambda
 \label{eq:irreg-chi}
\end{equation}
is nonzero. Under the Faddeev descent map, the invariant line is
generated by the class of \(1+p\).

Consider the following additional condition:
\begin{equation}
 \begin{gathered}
 d_n\in H^1(K,A)[\pi^n],\qquad
 \pi d_{n+1}=d_n \quad(n\geq1),\\
 d_{n,v}=0\quad(v\neq\mathfrak p),\qquad
 d_{1,\mathfrak p}\text{ pairs nontrivially with }\Lambda.
 \end{gathered}
 \label{eq:irreg-prescribed-tower}
\end{equation}
The dual of \(\pi\) is \(\overline\pi=1-\zeta_p^{-1}\), a unit
multiple of \(\pi\), so the last pairing is the usual finite local
Tate pairing on \(A(K_{\mathfrak p})/\pi A(K_{\mathfrak p})\).
Condition \eqref{eq:irreg-prescribed-tower} is stronger than the
mere existence of a nonzero supported \(p\)-primary tower.

\begin{proposition}\label{prop:irreg-prescribed-isolation}
If a wild quotient satisfies
\eqref{eq:irreg-prescribed-tower}, then any compact \(p\)-Selmer
datum on \(F_p\) whose \(p\)-local centre lies in a coordinate cusp
residue disc has that cusp as its centre.
\end{proposition}

\begin{proof}
The proof of Lemma~\ref{lem:irreg-reciprocity} applies over the
number field \(K\), with support at \(\mathfrak p\). Restrict a
compact datum to \(K\), push it to \(A\), and apply that argument
to every \(d_n\), using a finite \(p\)-power level killing \(d_n\).
The point \(j(f(R_p))\in A(\Q_p)\) lies in the common annihilator
\[
 B_K=\{P\in A(K_{\mathfrak p}):
             \langle P,d_{n,\mathfrak p}\rangle=0
             \text{ for every }n\}.
\]

Compatibility of the tower implies
\[
 \langle\overline\pi P,d_{n+1,\mathfrak p}\rangle
       =\langle P,d_{n,\mathfrak p}\rangle.
\]
Hence \(B_K\) is stable under \(\overline\pi\) and saturated for
multiplication by \(\overline\pi\). It is stable under
\(\Z[\zeta_p]\), and \(\overline\pi^{p-1}\) is \(p\) times a unit
of this ring. It is therefore saturated for multiplication by
\(p\). Multiplication by an integer prime to \(p\) is invertible
on every pairing target, giving saturation for all nonzero
integers. The same compatibility shows that \(B_K\) contains
all torsion: use a sufficiently large power of \(\overline\pi\)
for the \(p\)-primary part, and the order for the prime-to-\(p\) part.

Put \(B=B_K\cap A(\Q_p)\). Since the last condition in
\eqref{eq:irreg-prescribed-tower} is a nonzero functional on the
one-dimensional space \(\Lambda\), one has \(B\subseteq\ker\chi\).
In particular \(\chi\) vanishes on torsion.
The quotient \(A(\Q_p)/A(\Q_p)_{\tors}\) is a free
\(\Z_p\)-module, and the image of \(B\) is a saturated submodule.
Consequently there is a saturated subgroup \(M\) containing \(B\),
of corank one modulo torsion, whose image modulo \(p\) is
\(\ker\chi\). To see this directly, choose a basis adapted to the
direct summand defined by \(B\), and lift the nonzero linear
functional \(\chi\) on the complementary quotient to a primitive
\(\Z_p\)-linear functional; take its kernel.

The integral local calculation of
\cite[Section 4, Propositions 6--7 and Theorem 4,
 pp.~577--580]{McCallum1994}
now applies to \(M\). It gives a nonzero holomorphic differential
on \(C\), annihilating \(M\) under the abelian logarithm, whose
reduction is nonzero at at least one of \(O_0,O_1,O_\infty\).
This part of the argument uses the wild model and the specified
reduction \(M\bmod p=\ker\chi\), not regularity of \(p\).
Its Coleman integral vanishes at every allowed centre and at the
three torsion cusp images.

At a cusp where the reduced differential is nonzero, the local
zero estimate in
\cite[Lemma 7, pp.~579--580]{McCallum1994}
shows that the integral has at most one zero in that residue disc.
Since the cusp is already a zero, it is the only one there.
The quotient map respects these coordinate residue classes
\cite[Lemma 8 and the proof of Theorem 1, p.~580]{McCallum1994},
and is injective on \(\Q_p\)-points by
Lemma~\ref{lem:irreg-quotient-free}. Thus one coordinate cusp disc
on \(F_p\) contains no other permitted centre.

Finally the rational automorphism
\[
 \tau(X:Y:Z)=(-Y:Z:X)
\]
cycles the three coordinate cusps. It carries compact data to
compact data: on cohomology, replace \(t\) by
\[
 \tau_*t+
 \kappa_p([\tau(Q_0)-Q_0]),
\]
and replace every \(R_v\) by \(\tau(R_v)\). Thus the same isolated
disc may be applied after either cyclic permutation. All three
coordinate cusp discs are isolated for every compact datum.
\end{proof}

The hypotheses \((\mathrm H1)\), \((\mathrm H2)\) in
\cite[p.~575]{McCallum1994} ask for supported classes at arbitrarily
large \(\pi\)-levels with the specified nontrivial pairing after
multiplication down to the first level. Condition
\eqref{eq:irreg-prescribed-tower} is a compatible version sufficient
for the compact Selmer argument above. A compatible lift of the
initial Kummer class corresponding to \(\zeta_p\) is one sufficient
way to construct it; equivalence with that particular initial
class is not asserted.

\subsection{The obstruction space does not decide the specified lift}

Let \(\mathcal U=\Spec\mathcal O_K[1/p]\). Over \(K\), choose the
usual isomorphism \(A[\pi]\simeq\mu_p\). The Kummer sequence gives
\begin{equation}
 H^2(\mathcal U,A[\pi])
 \simeq H^2(\mathcal U,\mu_p)
 \simeq\Cl(K)/p\Cl(K).
 \label{eq:irreg-obstruction-space}
\end{equation}
Here \(\mathfrak p\) is principal, so
\(\operatorname{Pic}(\mathcal U)=\Cl(K)\). Also
\(\operatorname{Br}(\mathcal U)=0\): a class in this Brauer group
has zero invariant away from \(\mathfrak p\), all infinite places
of \(K\) are complex, and global reciprocity forces its remaining
invariant to vanish. These are the calculations in
\cite[Section 5, Theorem 5, pp.~581--582]{McCallum1994} and
\cite[Lemma 6, pp.~401--402]{McCallum2001Brauer}.

The regular-prime lifting construction uses the vanishing of
\eqref{eq:irreg-obstruction-space} to obtain the successive
surjections in cohomology associated with
\begin{equation}
 0\longrightarrow A[\pi]\longrightarrow A[\pi^{n+1}]
   \xrightarrow{\pi}A[\pi^n]\longrightarrow0.
 \label{eq:irreg-lift-sequence}
\end{equation}
At \(p=37,59,67\), the group in
\eqref{eq:irreg-obstruction-space} has dimension one.
This identifies a possible obstruction space. It does not compute
the obstruction to lifting the specified initial class, and does
not prove that such a lift fails.

In fact the first local pairing direction itself is present:
\cite[Lemma 10, pp.~581--582]{McCallum1994}
identifies the initial class with \(\zeta_p\) and its pairing with
the invariant line with the nonzero Hilbert symbol of
\(\zeta_p\) and \(1+p\).
The unresolved issue in this approach is a compatible lift through
arbitrarily high levels preserving the required support and
pairing direction.

Proposition~\ref{prop:irreg-McCallum}(1) constructs some nonzero
supported tower under \eqref{eq:irreg-Mc-bound}. Its statement
does not prescribe the first-layer pairing in
\eqref{eq:irreg-prescribed-tower}, so it cannot be substituted for
that stronger hypothesis. Likewise, the final observation in
\cite[pp.~405--406]{McCallum2001Brauer} that an individual complex
multiplication eigendifferential has only the three cusp zeros
does not show that such an eigendifferential annihilates all the
Selmer centres. The differential used in
Theorem~\ref{thm:irreg-fixed} is not asserted to have that property.

\subsection{Failure of the regular local-source surjectivity}

\begin{proposition}\label{prop:irreg-source-failure}
Suppose the \(p\)-primary part of \(\Cl(K)\) is nonzero. Let
\(K_{\{p\}}(p)\) be the maximal pro-\(p\) extension of \(K\)
unramified outside \(\mathfrak p\). Then the decomposition map
\begin{equation}
 G_{K_{\mathfrak p}}\longrightarrow
 \Gal(K_{\{p\}}(p)/K)
 \label{eq:irreg-source-map}
\end{equation}
is not surjective.
\end{proposition}

\begin{proof}
The maximal unramified abelian \(p\)-extension \(H/K\) is nontrivial
by class field theory and is contained in \(K_{\{p\}}(p)\).
The prime \(\mathfrak p=(1-\zeta_p)\) is principal, so its Artin
class in \(\Gal(H/K)\) is trivial. Since \(H/K\) is unramified,
\(\mathfrak p\) splits completely in \(H\). Its decomposition
group therefore maps trivially to the nontrivial quotient
\(\Gal(H/K)\), proving the assertion.
\end{proof}

For \(p=37,59,67\), the maximal unramified abelian \(p\)-extension
has degree \(p\). It is stable under \(\Delta\) and hence Galois
over \(\Q\); a decomposition group at \(p\) in \(\Gal(H/\Q)\) has
order \(p-1\), while the full group has order \(p(p-1)\).
Retaining the full arithmetic quotient therefore does not restore
this particular surjectivity input.

Proposition~\ref{prop:irreg-source-failure} is a statement about the
Galois source. It neither constructs a kernel for localization on
a particular geometric Tate module nor disproves an integral
restriction theorem for that module. A replacement restriction
argument, if needed for a global classification, must be supplied
separately. Even exact identification of all proper local centres
would not by itself classify sections in the remaining affine
geometric kernel.

